% =====================================================================
%  chunk c02 —— 植田一石《ホモロジー的ミラー対称性》
%  第 4–6 节，原书 p010–p013
%  编号公式 22 条（p010: 7 条，p011: 6 条，p012: 7 条，p013: 2 条），
%  逐页据 work/pages/pNNN.png 抄录，符号一一对应。
%  注：preamble 未载入 amscd / tikz-cd，原文的交换图表一律用 array 环境排出；
%  竖直箭头的标号按原版式置于箭头左侧或右侧。
% =====================================================================
% [Tutor Note: 关于本 chunk 的公式条数。work/BRIEF.md 对 c02 记「36 条编号公式」，
%  但按页面图像逐页核对，p010–p013 上实际印出的编号公式共 22 条，即
%  (4.1)–(4.19)、(5.1)、(6.1)、(6.2)。BRIEF 的 36 来自 work/tags_bypage.txt 的
%  原始命中数（p010: 8，p011: 9，p012: 17，p013: 2，合计 36），其中包含正文中的
%  公式交叉引用（如 p012 正文反复引用 (4.16)、(4.19)）与同一 tag 的重复命中，
%  并非编号公式的条数。此处按图像一一对应抄录 22 条（每页 7/6/7/2 条），
%  未增补、未重复、未自造任何编号。按 Lead 通知：以页面图像实际条数为准。]

……，右正合性、左正合性以及正合性是等价的概念。以稳定 ∞-范畴为对象、以正合函子为态射的 $\mathcal{C}at_\infty$ 的子范畴记作 $\mathcal{C}at^{\mathrm{Ex}}_\infty$。对特征 $0$ 的域 $k$，可以自然地定义 $k$ 上的稳定 ∞-范畴这一概念。$k$ 上的稳定 ∞-范畴所构成的 ∞-范畴，与把 $k$ 上的 dg 范畴所构成的范畴按森田等价作局部化而得到的 ∞-范畴等价。\footnote{参见 [Coh]。}

\section{A∞ 代数}

本节中，为了忘掉符号与分次的问题，我们在特征 $2$ 的域 $k$ 上考虑。本节的目标是解说弱 A∞ 代数与张量余代数上的余微分是等价的概念。\footnote{张量余代数上的余微分可以解释为非交换形式概形上的奇向量场（odd vector field）。关于这一点，参见 [KS09] 及其参考文献。}

设 $A$ 为 $k$ 上的向量空间，$m\colon A\otimes A\to A$ 与 $e\colon k\to A$ 为线性映射。所谓 $k$ 代数（k-algebra）指的是这样的 3 元组 $(A,m,e)$：图表
\begin{equation}
  \begin{array}{ccc}
    A\otimes A\otimes A & \xrightarrow{\ \id\otimes m\ } & A\otimes A\\[8pt]
    {\scriptstyle m\otimes\id}\Big\downarrow & & \Big\downarrow{\scriptstyle m}\\[8pt]
    A\otimes A & \xrightarrow{\ m\ } & A
  \end{array} \tag{4.1}
\end{equation}
可交换，并且
\begin{equation}
  A\cong k\otimes A \xrightarrow{\ e\otimes\id\ } A\otimes A \xrightarrow{\ m\ } A \tag{4.2}
\end{equation}
与
\begin{equation}
  A\cong k\otimes A \xrightarrow{\ e\otimes\id\ } A\otimes A \xrightarrow{\ m\ } A \tag{4.3}
\end{equation}
成为恒等映射。此时，$m$ 称为 $A$ 的乘法（multiplication），$e$ 称为单位元（unit）。图表 (4.1) 的可交换性就是 $A$ 的结合律（associativity）。对 $k$ 代数 $(A,m,e)$ 与 $(A',m',e')$，使以下图表
\begin{equation}
  \begin{array}{ccc}
    A\otimes A & \xrightarrow{\ m\ } & A\\[8pt]
    {\scriptstyle\varphi\otimes\varphi}\Big\downarrow & & \Big\downarrow{\scriptstyle\varphi}\\[8pt]
    A'\otimes A' & \xrightarrow{\ m'\ } & A'
  \end{array} \tag{4.4}
\end{equation}
以及
\begin{equation}
  \begin{array}{ccc}
    k & \xrightarrow{\ e\ } & A\\[8pt]
    {\scriptstyle\id}\Big\downarrow & & \Big\downarrow{\scriptstyle\varphi}\\[8pt]
    k & \xrightarrow{\ e'\ } & A'
  \end{array} \tag{4.5}
\end{equation}
可交换的线性映射 $\varphi\colon A\to A'$ 称为从 $(A,m,e)$ 到 $(A',m',e')$ 的 $k$ 代数同态（homomorphism）。以 $k$ 代数为对象、以 $k$ 代数同态为态射的范畴称为 $k$ 代数范畴，记作 $\mathrm{alg}_k$。此外，使图表
\begin{equation}
  \begin{array}{ccc}
    A\otimes A & \xrightarrow{\ m\ } & A\\[8pt]
    {\scriptstyle\id\otimes\partial+\partial\otimes\id}\Big\downarrow & & \Big\downarrow{\scriptstyle\partial}\\[8pt]
    A\otimes A & \xrightarrow{\ m\ } & A
  \end{array} \tag{4.6}
\end{equation}
可交换的线性映射 $\partial\colon A\to A$ 称为 $A$ 上的导子（derivation）。

% [Tutor Note: (4.2) 与 (4.3) 在图像上完全相同，均为
%  A ≅ k⊗A —e⊗id→ A⊗A —m→ A。按常理 (4.3) 似应为
%  A ≅ A⊗k —id⊗e→ A⊗A —m→ A。此处照录图像，请审校确认。]

对 $k$ 代数 $A$，把其底层结构所对应的向量空间与之对应的遗忘函子（forgetful functor）记作 $\mathrm{fgt}\colon\mathrm{alg}_k\to\vec{t}_k$。它本质满（essentially surjective）且忠实（faithful），但不是满的（full）。对向量空间 $V$，在直和 $TV:=\bigoplus_{n\in\N}V^{\otimes n}$ 上按
\begin{equation}
  (v_1\otimes\cdots\otimes v_m)\otimes(w_1\otimes\cdots\otimes w_n)
  \mapsto v_1\otimes\cdots\otimes v_m\otimes w_1\otimes\cdots\otimes w_n \tag{4.7}
\end{equation}
引入乘法 $TV\otimes TV\to TV$ 而得到的代数，称为 $V$ 上的自由代数（free algebra），记作 $\mathrm{free}(V)$。\footnote{自由代数也称为张量代数（tensor algebra）。}

设 $\varphi\colon V\to V'$ 为从向量空间 $V$ 到 $V'$ 的线性映射，定义 $T\varphi\colon TV\to TV'$ 为 $T\varphi(v_1\otimes\cdots\otimes v_n):=\varphi(v_1)\otimes\cdots\otimes\varphi(v_n)$，则它成为 $k$ 代数同态。因此，把向量空间 $V$ 对应到张量代数 $TV$ 的函子 $\mathrm{free}\colon\vec{t}_k\to\mathrm{alg}_k$ 得以确定，并且它是遗忘函子 $\mathrm{fgt}$ 的左伴随函子（left adjoint functor）；
\begin{equation}
  \Hom_{\vec{t}_k}(V,\mathrm{fgt}(A)) \cong \Hom_{\mathrm{alg}_k}(\mathrm{free}(V),A). \tag{4.8}
\end{equation}
伴随关系 (4.8) 与下面的普遍映射性质（universal mapping property）等价：对任意线性映射 $f\colon V\to A$，存在唯一的代数同态 $\tilde f\colon\mathrm{free}(V)\to A$，使 $f$ 唯一地经由嵌入映射 $\varphi_V\colon V\hookrightarrow\mathrm{free}(V)$ 分解。

把代数的公理对偶化，便得到余代数的概念。$k$ 余代数（k-coalgebra）是指由向量空间 $C$、线性映射 $\Delta\colon C\to C\otimes C$ 及 $\epsilon\colon C\to k$ 组成的 3 元组 $(C,\Delta,\epsilon)$，满足图表
\begin{equation}
  \begin{array}{ccc}
    C & \xrightarrow{\ \Delta\ } & C\otimes C\\[8pt]
    {\scriptstyle\Delta}\Big\downarrow & & \Big\downarrow{\scriptstyle\id\otimes\Delta}\\[8pt]
    C\otimes C & \xrightarrow{\ \Delta\otimes\id\ } & C\otimes C\otimes C
  \end{array} \tag{4.9}
\end{equation}
可交换，并且
\begin{equation}
  C \xrightarrow{\ \Delta\ } C\otimes C \xrightarrow{\ \epsilon\otimes\id\ } k\otimes C \cong C \tag{4.10}
\end{equation}
及
\begin{equation}
  C \xrightarrow{\ \Delta\ } C\otimes C \xrightarrow{\ \id\otimes\epsilon\ } k\otimes C \cong C \tag{4.11}
\end{equation}
等于恒等映射这一条件。此时，$\Delta$ 称为余积（coproduct），$\epsilon$ 称为余单位元（counit）。此外，图表 (4.9) 的可交换性称为余结合律（coassociativity）。与代数的情形一样，与余积及余单位元可交换的线性映射称为 $k$ 余代数同态；以 $k$ 余代数为对象、以 $k$ 余代数同态为态射的范畴记作 $\mathrm{Coalg}_k$。对余代数 $C$，使图表
\begin{equation}
  \begin{array}{ccc}
    C & \xrightarrow{\ \Delta\ } & C\otimes C\\[8pt]
    {\scriptstyle\Delta}\Big\downarrow & & \Big\downarrow{\scriptstyle\id\otimes\delta+\delta\otimes\id}\\[8pt]
    C & \xrightarrow{\ \Delta\ } & C\otimes C
  \end{array} \tag{4.12}
\end{equation}
可交换的线性映射 $\delta\colon C\to C$ 称为 $C$ 上的余导子（coderivation）。此外，满足 $\delta\circ\delta=0$ 的余导子称为余微分（codifferential）。

% [Tutor Note: (4.11) 的像原文印作 k⊗C（箭头标号为 id⊗ε: C⊗C → k⊗C）；
%  按 (4.10) 的对称性，此处似应为 C⊗k。此处照录图像，请审校确认。]

对任意向量空间 $V$，直和 $TV:=\bigoplus_{n\in\N}V^{\otimes n}$ 带有以
\begin{equation}
  \Delta(v_1\otimes\cdots\otimes v_n) = \sum_{i=0}^{n}(v_1\otimes\cdots\otimes v_i)\otimes(v_{i+1}\otimes\cdots\otimes v_n) \tag{4.13}
\end{equation}
定义的映射 $\Delta\colon TV\to TV\otimes TV$ 为余积、以到直和分量 $k\cong V^{\otimes 0}\subset TV:=\bigoplus_{n\in\N}V^{\otimes n}$ 的投影 $\epsilon$ 为余单位元的余代数结构，这称为 $V$ 上的张量余代数（tensor coalgebra）。

若存在某个自然数 $n$ 使 $\Delta_n\colon C\to C^{\otimes n}$ 为 $0$，则称余代数 $C$ 是余幂零的（conilpotent）。可表示成余幂零余代数的归纳极限的余代数称为归纳余幂零的（ind-colinpotent）。把向量空间 $V$ 对应到 $V$ 上的张量余代数 $C(V)$ 的函子，是「从归纳余幂零余代数的范畴到向量空间的范畴的遗忘函子」的右伴随函子 [Pri10, Remark 3.12]\footnote{由自由代数 $\mathrm{free}(V)$ 的底层结构即作为向量空间由直和 $TV=\bigoplus_{n\in\N}V^{\otimes n}$ 给出可知，余自由代数 $\mathrm{cofree}(V)$ 的底层结构似乎取直积 $\widehat{TV}:=\prod_{n\in\N}V^{\otimes n}$ 即可；但张量余代数 $TV\subset\widehat{TV}$ 的余积 (4.13) 不能延拓为 $\widehat{TV}\to\widehat{T}\otimes\widehat{T}$，所以这种想当然的想法并不正确。关于余自由代数的构造，参见例如 [Haz03]。}。

设 $V$ 为向量空间，$\delta\colon TV\to TV$ 为张量余代数 $TV$ 上的余微分。记与直和分解 $TV:=\bigoplus_{i\in\N}V^{\otimes i}$ 相伴的投影为 $\pi\colon TV\to V$，并记复合映射 $\mathbf{m}:=\pi\circ\delta\colon TV\to V$ 在 $V^{\otimes n}\subset TV$ 上的限制为
\begin{equation}
  \mathbf{m}_n\colon V^{\otimes n}\to V \tag{4.14}
\end{equation}
则由图式 (4.12) 的可交换性，余导子 $\delta$ 由映射的序列 $(\mathbf{m}_n)_{n\in\N}$ 按
\begin{equation}
  \delta(v_1\otimes\cdots\otimes v_n)
  = \sum_{m=0}^{n}\sum_{i=0}^{n-m}
  v_1\otimes\cdots\otimes v_i\otimes \mathbf{m}_m(v_{i+1}\otimes\cdots\otimes v_{i+m})
  \otimes v_{i+m+1}\otimes\cdots\otimes v_n \tag{4.15}
\end{equation}
唯一地确定。这里，(4.15) 右端求和号内的项是 $V^{\otimes(n-m+1)}\subset TV$ 的元素。由此可知，余导子 $\delta$ 是余微分的充要条件是：对任意 $n\in\N$ 与任意 $v_1\otimes\cdots\otimes v_n\in V^{\otimes n}$，
\begin{equation}
  \sum_{m=0}^{n}\sum_{i=0}^{n-m}
  \mathbf{m}_{n-m+1}\bigl(v_1\otimes\cdots\otimes v_i\otimes \mathbf{m}_m(v_{i+1}\otimes\cdots\otimes v_{i+m})
  \otimes v_{i+m+1}\otimes\cdots\otimes v_n\bigr) = 0 \tag{4.16}
\end{equation}
成立。(4.16) 称为 A∞ 关系式（A∞-relation）；由向量空间 $V$ 与映射的序列 $(\mathbf{m}_n)_{n\in\N}$ 组成的组 $(V,(\mathbf{m}_n)_{n\in\N})$ 中满足 (4.16) 者，称为弱 A∞ 代数（weak A∞-algebra）。弱 A∞ 代数进一步满足条件 $\mathbf{m}_0=0$ 时，称为 A∞ 代数（A∞-algebra）。在 A∞ 代数的情形，把较小的 $n$ 所对应的 A∞ 关系式 (4.16) 具体写出来，就是
\begin{equation}
  \mathbf{m}_1(\mathbf{m}_1(x)) = 0, \tag{4.17}
\end{equation}
\begin{equation}
  \mathbf{m}_2(\mathbf{m}_1(x)\otimes y) + \mathbf{m}_2(x\otimes\mathbf{m}_1(y)) + \mathbf{m}_1(\mathbf{m}_2(x\otimes y)) = 0, \tag{4.18}
\end{equation}
\begin{align}
  & \mathbf{m}_3(\mathbf{m}_1(x)\otimes y\otimes z) + \mathbf{m}_3(x\otimes\mathbf{m}_1(y)\otimes z)
    + \mathbf{m}_3(x\otimes y\otimes\mathbf{m}_1(z)) \notag\\
  & \quad + \mathbf{m}_2(\mathbf{m}_2(x\otimes y)\otimes z) + \mathbf{m}_2(x\otimes\mathbf{m}_2(y\otimes z))
    + \mathbf{m}_1(\mathbf{m}_3(x\otimes y\otimes z)) = 0 \tag{4.19}
\end{align}
成为这样。(4.17) 意味着 $\mathbf{m}_1$ 是微分，(4.18) 是乘法 $\mathbf{m}_2$ 关于微分 $\mathbf{m}_1$ 的 Leibniz 法则。(4.19) 在 $\mathbf{m}_3=0$ 时给出 $\mathbf{m}_2$ 的结合律。当 $\mathbf{m}_3\neq 0$ 时 $\mathbf{m}_2$ 不是结合的，因此一般地 A∞ 范畴不成为范畴\footnote{一般的 A∞ 范畴不是范畴，这与任意微分分次范畴都是范畴形成对照。}。但是，由 Leibniz 法则 (4.18)，$\mathbf{m}_2$ 诱导出上同调群 $H^{*}(V,\mathbf{m}_1):=\operatorname{Ker}\mathbf{m}_1/\operatorname{Im}\mathbf{m}_1$ 上的乘法 $[\mathbf{m}_2]\colon H^{*}(V,\mathbf{m}_1)\otimes H^{*}(V,\mathbf{m}_1)\to H^{*}(V,\mathbf{m}_1)$，而 (4.19) 保证这个乘法是结合的。换言之，$\mathbf{m}_2$ 在相差同伦的意义下是结合的，$\mathbf{m}_3$ 描述的正是这个同伦。进而，$\mathbf{m}_3$ 在相差同伦的意义下满足高阶的相容性条件，其同伦由 $\mathbf{m}_4$ 描述，……，这样的链条无限延续下去。

\section{导出代数几何}

增广三角范畴是加性的（或者线性的）理论，但也可以考虑非线性对象——代数流形——的导出版本。其出发点是把代数流形看作从可换环的范畴 $\mathrm{calg}$ 到集合的范畴 $\mathrm{set}$ 的函子。从给定的代数流形造出新的代数流形的典型方法，有切（即取函数的零点）这一操作与除（即取按等价关系的商）这一操作。函数 $f\colon X\to\mathbb{A}^1$ 的零点就是图表
\begin{equation}
  \begin{array}{ccc}
    & & X\\[8pt]
    & & \Big\downarrow{\scriptstyle f}\\[8pt]
    0 & \longrightarrow & \mathbb{A}^1
  \end{array} \tag{5.1}
\end{equation}
的纤维积，而按等价关系取的商则是图表 $S\overset{f}{\underset{g}{\rightrightarrows}}X$ 的余等化子（coequalizer）。前者在坐标环的层面上对应于张量积，而张量积不是左正合的这一点，在非横截（transversal）的情形会引起问题。后者的典型例子是群作用下的商，它往往无法用概形表示。前者的困难可以通过把函子的定义域从 $\mathrm{calg}$ 换成单纯可换环（simplicial commutative ring）的范畴 $\mathrm{scalg}$ 来消解。另一方面，把函子的值域换成群胚（groupoid），从而使按群作用的商总是可以被表示出来，这便是栈（stack）的理论。但是，为了处理比余等化子更一般的余极限（colimit），考虑在高阶范畴中与群胚相对应的 ∞-群胚（∞-groupoid）会更方便。把 Grothendieck 的同伦假设（homotopy hypothesis）记在心上，取单纯集合的范畴 $\mathrm{sset}$ 作为 ∞-群胚的模型，并把从 $\mathrm{scalg}$ 到 $\mathrm{sset}$ 的、满足适当条件的函子定义为导出代数流形。这样得到的导出代数流形的范畴，提供了可以自由地取极限与余极限的灵活框架。关于导出代数几何的概说，参见例如 [Toë09, Toë14]。

作为有代表性的导出代数流形，可以举出代数流形 $X$ 的环路空间 $LX:=\operatorname{Map}(S^1,X)=X\times_{X\times X}X$。它是把 $\cO_X$ 的 Hochschild 复形看作 $\cO_X$ 上的微分分次代数时的谱。由 $S^1$ 等变的 Hochschild 上同调是循环上同调，以及代数流形的循环上同调是 de Rham 上同调这两点可知，代数流形上的 $\cD$-模可以用其环路空间上 $S^1$ 等变的准相干层来描述 [BZN12, TV15]。此外，栈的环路空间也同样被定义，可以看作惯性栈（inertia stack）的导出版本。

对于只对经典代数流形感兴趣的人来说，导出代数几何就如同对于只对代数方程的整数解感兴趣的人来说的复代数几何。即使在结论中并不出现导出代数流形，为了自然地理解而有效运用导出代数几何的例子，可以举出它在几何表示论中的应用 [BZN13]。

在镜像对称中，各种模空间起着中心作用。模空间常常是通过把方程的零点按等价关系作商而得到的，因此用导出代数几何的方式来处理是自然的。把模空间构造成导出栈，从而消解「在代数流形或概形的范围内考虑模问题时产生的各种困难」——这一哲学称为隐光滑性原理（hidden smoothness principle）\footnote{藤原一宏，《现代数学的土壤》伽罗瓦理论，数学的乐趣，第 13 卷，日本评论社，pp. 109–121 中把 Galois 理论比作隐对称性（hidden symmetry），而隐光滑性也与隐对称性略有相似。隐对称性也称为自发破缺的对称性（spontaneously broken symmetry），是理论物理学中的重要原理。要直接看到隐对称性需要高能量，而要直接看到隐光滑性则需要高阶范畴。}。此外，用微分 Lie 环描述模空间的局部理论，也能在导出代数几何的框架中自然地理解 [Man02, Hin01, Lur10]。

% [Tutor Note: (5.1) 的图表中，方图顶边 X → 0 的箭头未印出，图上只画出 0 → A^1 与
%  X —f→ A^1 两个箭头；此处照录图像。该图表表示 X 是 0 → A^1 沿 f: X → A^1 的纤维积。
%  另，图注「図式 S ⇉ X」中上下两个箭头分别标 f、g，已按原版式排出。]

\section{复几何与辛几何}

给定 $m$ 维流形 $M$、群 $G$，以及从 $G$ 到 $\mathrm{GL}(m,\R)$ 的同态时，$M$ 的标架丛 $TM$ 的结构群 $\mathrm{GL}(m,\R)$ 到 $G$ 的约化称为 $M$ 的 $G$ 结构（G-structure）。例如，$\mathrm{GL}(m,\R)^{+}$ 结构与定向等价，$O(m)$ 结构与黎曼度量等价。此外，当给定的 $G$ 结构局部同构于 $\R^m$ 的自然 $G$ 结构时，称为可积的（integrable）。定向总是可积的，而黎曼结构可积的充要条件是曲率张量为 $0$。

复一般线性群（complex general linear group）与辛群（symplectic group）用矩阵 $J_{2n}:=\begin{pmatrix}0 & -I_n\\ I_n & 0\end{pmatrix}$ 定义为
\begin{equation}
  \mathrm{GL}(n,\C) := \{\, g\in\mathrm{GL}(2n,\R) \mid g^{-1}J_{2n}g = J_{2n} \,\}, \tag{6.1}
\end{equation}
\begin{equation}
  \mathrm{Sp}(2n,\R) := \{\, g\in\mathrm{GL}(2n,\R) \mid g^{T}J_{2n}g = J_{2n} \,\} \tag{6.2}
\end{equation}
可积的 $\mathrm{Sp}(2n,\R)$ 结构称为辛结构（symplectic structure），可积的 $\mathrm{GL}(n,\C)$ 结构称为复结构（complex structure）。在偶数维流形上给出 $\mathrm{Sp}(2n,\R)$ 结构，等价于给出在所有点处非退化的 2 次微分形式 $\omega$；而给出 $\mathrm{GL}(n,\C)$ 结构，等价于给出切丛 $TM$ 的自同态 $J\in\End(TM)$ 且满足 $J^2=-\id_{TM}$\footnote{未必可积的 $\mathrm{GL}(n,\C)$ 结构称为近复结构（almost complex structure）。未必可积的 $\mathrm{Sp}(2n,\R)$ 结构没有特别的称呼。}。满足 $Jf_{*}=f_{*}J$ 的微分同胚称为双全纯映射（biholomorphic map），满足 $f^{*}\omega=\omega$ 的微分同胚称为辛同胚（symplectomorphism）。此外，辛向量空间的开集之间的辛同胚特别地称为正则变换（canonical transformation）。复流形是用双全纯映射把复向量空间的开集粘贴起来所得的流形，而辛流形则是把辛向量空间的开集用正则变换

% CHUNK c02 DONE pages p010–p013
